% ============================================================================
% CAPÍTULO 9 — Estilização moderna e design systems
% Texto completo (rodada editorial 2)
% ============================================================================
\chapter{Estilização Moderna e Design Systems}
\label{cap:estilizacao}

\begin{caixaconceito}
\textbf{Tailwind CSS v4} \cite{tailwindv4} é o framework de estilos
\emph{utility-first} mais usado do mercado: em vez de escrever classes
semânticas e CSS separado, você compõe a interface com utilitários atômicos
(\texttt{flex}, \texttt{p-4}, \texttt{text-center}) — e o v4 faz isso com
configuração CSS-first, sem arquivo de configuração JavaScript.
\end{caixaconceito}

Ao final deste capítulo, você será capaz de:
\begin{objetivos}
  \item Comparar as abordagens de estilização (CSS Modules, CSS-in-JS, utility-first\index{utility-first});
  \item Configurar o Tailwind v4 em um projeto (Vite e Next.js);
  \item Aplicar utilitários, responsividade, dark mode\index{dark mode} e variantes;
  \item Criar um design system\index{design system} com tokens e componentes reutilizáveis;
  \item Integrar componentes de UI (shadcn/ui\index{shadcn/ui}) sem perder o controle;
  \item Entregar o projeto do capítulo: o design system \emph{OrçaUI}.
\end{objetivos}

\section{As abordagens de estilização no mercado}

Em 2026, convivem quatro abordagens — e é importante saber \emph{onde} cada
uma se encaixa:

\begin{itemize}[leftmargin=*, itemsep=0.4em]
  \item \textbf{CSS puro/Módulos}: você já domina a base (capítulo~\ref{cap:css});
  CSS Modules adiciona escopo por componente (\texttt{Button.module.css});
  \item \textbf{CSS-in-JS} (styled-components, Emotion): estilos como
  JavaScript, com temas dinâmicos — caiu em uso por custo de runtime e
  problemas com Server Components;
  \item \textbf{Utility-first} (Tailwind): classes atômicas no HTML;
  \emph{padrão dominante} em 2026;
  \item \textbf{Design systems prontos} (shadcn/ui, Radix, MUI): componentes
  acessíveis e customizáveis por cima do Tailwind.
\end{itemize}

\begin{caixaconceito}
O debate ``Tailwind vs CSS puro'' é falso: você \emph{precisa} entender CSS
(capítulo~\ref{cap:css}) para usar Tailwind bem. Tailwind é uma camada de
produtividade sobre o mesmo modelo de cascata, especificidade e box model —
as mesmas regras que você aprendeu, com nomes curtos.
\end{caixaconceito}

\section{Tailwind v4: configuração CSS-first}

No v4, o design system vive em CSS com a diretiva \texttt{@theme} — sem
\texttt{tailwind.config.js}:

\begin{lstlisting}[style=livro, language=CSS,
  caption={Tokens do design system no Tailwind v4.}, label={list:tailwind-tema}]
@import "tailwindcss";

@theme {
  --color-primaria: #1f4e79;
  --color-secundaria: #2e86ab;
  --color-acento: #f4a261;
  --color-fundo: #f7f9fb;
  --color-texto: #22313f;

  --font-display: "Inter", system-ui, sans-serif;
  --font-corpo: "Inter", system-ui, sans-serif;

  --radius-card: 12px;
  --shadow-suave: 0 4px 12px rgb(31 78 121 / 0.12);
}
\end{lstlisting}

\begin{lstlisting}[style=livro, language=Bash,
  caption={Instalando o Tailwind v4 no Vite.}, label={list:tailwind-instalar}]
npm install tailwindcss @tailwindcss/vite
# vite.config.ts: adicione tailwindcss() aos plugins
# e no CSS principal: @import "tailwindcss";
\end{lstlisting}

\begin{caixaconceito}
Os tokens definidos em \texttt{@theme} viram \emph{utilitários}:
\texttt{bg-(--color-primaria)} (v4) ou \texttt{bg-[var(--color-primaria)]}
(v3). No v4, você também pode gerar utilitários nomeados automaticamente —
por exemplo, \texttt{--color-primaria} cria \texttt{text-primaria},
\texttt{bg-primaria}, \texttt{border-primaria}.
\end{caixaconceito}

\section{Utilitários, responsividade e dark mode}

\begin{lstlisting}[style=livro, language=HTML,
  caption={Composição com utilitários Tailwind.}, label={list:tailwind-uso}]
<button
  class="rounded-lg bg-primaria px-6 py-3 font-bold text-white
         transition hover:bg-secundaria
         focus-visible:outline focus-visible:outline-2
         md:px-8 dark:bg-gray-800"
>
  Começar agora
</button>
\end{lstlisting}

\begin{itemize}[leftmargin=*, itemsep=0.4em]
  \item \textbf{Responsividade}: prefixos \texttt{sm:}, \texttt{md:},
  \texttt{lg:}, \texttt{xl:} aplicam o utilitário a partir da largura
  (mobile-first, como no capítulo~\ref{cap:css});
  \item \textbf{Dark mode}: \texttt{dark:} ativa quando a classe \texttt{dark}
  está no \texttt{<html>} (ou com \texttt{prefers-color-scheme});
  \item \textbf{Variantes de estado}: \texttt{hover:}, \texttt{focus-visible:},
  \texttt{active:}, \texttt{disabled:}, \texttt{group-hover:}, \texttt{first:}, \texttt{even:};
  \item \textbf{Valores arbitrários} (exceção, não regra):
  \texttt{bg-[\#123456]}, \texttt{w-[17px]}.
\end{itemize}

\begin{caixadica}
Regra de ouro de utilitários: \textbf{padrão primeiro}. Use as escalas padrão
(\texttt{p-4}, \texttt{text-lg}) sempre que possível; valores arbitrários são
exceção. Isso mantém o design consistente e o bundle pequeno — e facilita
mudanças globais: altere o token, e todos os componentes seguem.
\end{caixadica}

\section{Extraindo componentes: do utilitário ao design system}

Utilitários no HTML ficam repetitivos quando o mesmo botão aparece 20 vezes.
A solução é \emph{extrair} para componentes — e é aí que o React brilha:

\begin{lstlisting}[style=livro, language=TypeScript,
  caption={Componente Button com variantes.}, label={list:tailwind-componente}]
import { type ButtonHTMLAttributes } from "react";

type Variante = "primario" | "secundario" | "perigo";

const estilos: Record<Variante, string> = {
  primario: "bg-primaria text-white hover:bg-secundaria",
  secundario: "border-2 border-primaria text-primaria hover:bg-primaria hover:text-white",
  perigo: "bg-red-600 text-white hover:bg-red-700",
};

interface BotaoProps extends ButtonHTMLAttributes<HTMLButtonElement> {
  variante?: Variante;
}

export function Botao({ variante = "primario", className = "", ...props }: BotaoProps) {
  return (
    <button
      className={`rounded-lg px-6 py-3 font-bold transition
        focus-visible:outline focus-visible:outline-2 ${estilos[variante]} ${className}`}
      {...props}
    />
  );
}
\end{lstlisting}

\begin{caixaconceito}
O padrão \emph{componente + variantes} é o coração dos design systems:
\emph{um} componente \texttt{Botao} com 3 variantes substitui 20 marcações
repetidas. Se amanhã o ``raio'' muda, você altera um lugar. É a ponte entre
``utility-first no HTML'' e ``design system no produto''.
\end{caixaconceito}

\section{Design systems: tokens, componentes e documentação}

Um design system é a ``língua visual'' do produto, com três camadas:

\begin{enumerate}[leftmargin=*, itemsep=0.4em]
  \item \textbf{Tokens}: cores, espaçamentos, raios, tipografia, sombras
  (definidos em \texttt{@theme});
  \item \textbf{Componentes base}: \texttt{Button}, \texttt{Input},
  \texttt{Card}, \texttt{Badge} com variantes e estados;
  \item \textbf{Padrões de página}: como os componentes se combinam
  (formulários, tabelas, navegação) — com exemplos vivos.
\end{enumerate}

\begin{caixaconceito}
\textbf{shadcn/ui} é a abordagem dominante em 2026: em vez de uma biblioteca
fechada, você \emph{copia} componentes acessíveis (construídos sobre Radix
UI) para o seu projeto e os customiza livremente. O código é seu — o que casa
perfeitamente com o modelo deste livro e evita a ``dívida'' de bibliotecas
imutáveis.
\end{caixaconceito}

\section{Acessibilidade visual: contraste e foco}

\begin{itemize}[leftmargin=*, itemsep=0.4em]
  \item \textbf{Contraste}: texto sobre fundo $\geq$ 4,5:1 (WCAG\index{WCAG} AA) — valide
  com a ferramenta do Lighthouse (capítulo~\ref{cap:performance});
  \item \textbf{Foco visível}: sempre \texttt{focus-visible:outline-2} nos
  interativos — nunca remova o outline sem substituto;
  \item \textbf{Reduced motion}: respeite \texttt{prefers-reduced-motion} nas
  animações;
  \item \textbf{Tamanho de toque}: alvos de interação $\geq$ 44px em mobile.
\end{itemize}

% ============================================================================
\section{Projeto do capítulo: o design system OrçaUI}
\label{sec:projeto9}

\begin{caixaprojeto}
\textbf{Objetivo:} criar o design system \emph{OrçaUI} para o produto fictício
OrçaFácil (capítulo~\ref{cap:css}): tokens, componentes base e uma página de
documentação (``storybook caseiro'').

\textbf{Critérios de aceite:}
\begin{itemize}[leftmargin=*, itemsep=0.2em]
  \item Tokens completos em \texttt{@theme} (cores, tipografia, espaçamentos, raios, sombras);
  \item Componentes \texttt{Button} (4 variantes), \texttt{Input}, \texttt{Card}, \texttt{Badge}, \texttt{Modal} — com TypeScript;
  \item Estados: hover, foco visível, desabilitado e carregando;
  \item Dark mode funcional em todos os componentes;
  \item Página \texttt{/design-system} documentando cada componente com exemplos vivos;
  \item Testes visuais básicos (screenshot) — configurados com Playwright no capítulo~\ref{cap:testes}.
\end{itemize}
\end{caixaprojeto}

\begin{caixadica}
Um \emph{storybook caseiro} é uma página que mostra cada componente em todos
os seus estados (padrão, hover, erro, carregando, dark). É a forma mais
rápida de ``ver'' o design system crescer — e o mesmo conceito das
ferramentas profissionais (Storybook, Chromatic).
\end{caixadica}

\section{Exercícios de fixação}

\begin{exercicios}
  \item Compare CSS Modules, CSS-in-JS e utility-first: uma vantagem e uma desvantagem de cada.
  \item Escreva a diretiva \texttt{@theme} com 3 tokens de cor e 2 de raio.
  \item Crie um botão responsivo com Tailwind: centralizado, com padding em \texttt{md:} maior e cor de fundo no hover.
  \item O que são variantes \texttt{group-hover} e \texttt{focus-visible}? Dê um uso real.
  \item Liste as três camadas de um design system e dê um exemplo de cada.
  \item Por que extrair componentes com variantes em vez de repetir utilitários?
\end{exercicios}

\section{Desafios}

\begin{desafios}
  \item \textbf{Dark mode completo}: implemente o modo escuro com botão de alternância e persistência em \texttt{localStorage} em um app com 3 páginas.
  \item \textbf{Componente Modal acessível}: construa um \texttt{Modal} com foco preso (focus trap), fechamento por \texttt{Escape} e \texttt{role="dialog"} + \texttt{aria-modal}.
  \item \textbf{Tokens responsivos}: use \texttt{clamp()} nos tokens de tipografia para que títulos escalem suavemente entre 320px e 1440px.
  \item \textbf{Badge e variantes}: implemente o \texttt{Badge} com 4 tons (info, sucesso, alerta, perigo) e teste o contraste de cada um.
\end{desafios}

\section{Exercícios de nível sênior}

\begin{seniores}
  \item \textbf{Anatomia do bundle}: analise com \texttt{npx next build} (ou Vite) quanto CSS cada página carrega e otimize removendo utilitários redundantes.
  \item \textbf{Contraste e tokens}: audite os tokens de cor do OrçaUI contra WCAG AA e ajuste a paleta (documente a decisão de cada cor).
  \item \textbf{Componente composto}: implemente um \texttt{Tabs} com compound components (\texttt{Tabs.List}, \texttt{Tabs.Panel}) seguindo o padrão do Radix UI, com navegação por teclado (setas).
  \item \textbf{Tema no Next.js}: implemente tokens dinâmicos no App Router com variáveis CSS por tema e \texttt{class} no \texttt{<html>}, sem flash de tema incorreto (FOUC).
\end{seniores}

\section{Armadilhas comuns}

\begin{itemize}[leftmargin=*, itemsep=0.4em]
  \item Usar valores arbitrários em excesso (design vira ``sopa de pixels'');
  \item Aplicar Tailwind sem entender CSS (resultado: layouts quebrados);
  \item Dark mode ``na gambiarra'' com cores fixas fora de tokens;
  \item Copiar bibliotecas de UI inteiras sem questionar acessibilidade e tamanho;
  \item Esquecer \texttt{focus-visible} — o teclado fica invisível.
\end{itemize}

\section{Resumo do capítulo}

\begin{itemize}[leftmargin=*, itemsep=0.3em]
  \item Utility-first (Tailwind v4) é o padrão de estilização de 2026;
  \item \texttt{@theme} centraliza os tokens sem arquivo JS de configuração;
  \item Responsividade (\texttt{md:}), dark mode (\texttt{dark:}) e variantes são prefixos;
  \item Componentes com variantes substituem marcações repetidas;
  \item Design systems = tokens + componentes base + padrões de página;
  \item shadcn/ui entrega acessibilidade com código que é seu;
  \item Projeto entregue: design system OrçaUI documentado.
\end{itemize}

\section{Referências oficiais para aprofundar}

\begin{itemize}[leftmargin=*, itemsep=0.3em]
  \item Documentação oficial do Tailwind: \url{https://tailwindcss.com/docs}
  \item Anúncio do Tailwind v4: \url{https://tailwindcss.com/blog/tailwindcss-v4}
  \item shadcn/ui: \url{https://ui.shadcn.com}
  \item Radix UI (primitivos acessíveis): \url{https://www.radix-ui.com}
  \item WCAG 2.2 (W3C): \url{https://www.w3.org/WAI/WCAG22/Overview}
\end{itemize}
